\chapter{Technologie i narzędzia}
\label{chap:technologie}
Wybrane technologie i narzędzia miały za zadanie zrealizować wymagania projektu jak wieloplatformowość, niezależność od CUDA, czas rzeczywisty dla scen rzędu milionów Gaussianów, jak i stworzyć warstwę interfejsu i analizy. Takie wymagania sprawiają, że nie jesteśmy uzależnieni od kart NVIDIA, zostają wykluczone języki interpretowane w warstwie renderującej i musi istnieć GUI potrzebne do analizy.
\section{Język C++ (standard C++20)}
\label{sec:c++}
C++ to wszechstronny kompilowany język programowania ogólnego przeznaczenia. Powstał jako rozszerzenie języka C o mechanizmy programowania obiektowego, umożliwia abstrakcję danych oraz stosowanie kilku paradygmatów programowania: proceduralnego, obiektowego, generycznego, modularnego i funkcjonalnego. Język ten został wybrany ze względu na ręczne zarządzanie pamięcią i zerowy narzut abstrakcji. Wizualizacja scen 3DGS wymaga przetwarzania milionów splatów w każdej klatce, a każda warstwa pomiędzy kodem a sterownikiem graficznym przekłada się na wolniejsze wykonywanie kodu źródłowego. C++ zapewnia pełną kontrolę nad układem danych w pamięci, co pomaga przygotować dane do wysłania na GPU bez kosztownych konwersji. Z mechanizmów standardu C++20 wykorzystano \texttt{std::bit\_cast} z nagłówka \texttt{<bit>}, użyty przy budowie kluczy sortowania po głębokości~\cite{wg21_2020draft}. Porównywane jest sortowanie na GPU z sortowaniem równoległym na CPU (\texttt{std::execution::par\_unseq}, mechanizm wprowadzony w C++17). Język ten jest popularny w środowisku grafiki komputerowej -- API graficzne i biblioteki matematyczne udostępniają interfejsy w C lub C++.
\section{OpenGL 4.3 Core Profile}
\label{sec: opengl}
OpenGL to otwarta, wieloplatformowa i uniwersalna specyfikacja do tworzenia grafiki 3D, rozwijana przez Khronos Group. Podstawową cechą OpenGL jest to, że jest on tzw. maszyną stanu \eng{state machine}. Na stan w danej chwili składają się parametry i tryby działania ustawiane wywołaniami funkcji; raz ustawiona wartość obowiązuje aż do jej zmiany i wpływa na wynik renderowania. W projekcie wykorzystano wersję 4.3 Core, co pozwala wykorzystać shadery obliczeniowe \eng{compute shaders} oraz Shader Storage Buffer Object, czyli bufory pozwalające na zapis i odczyt dużych struktur danych z poziomu shadera~\cite{khronos2013opengl43}. Te dwa mechanizmy pozwalają zaimplementować potok renderowania 3DGS bez użycia CUDA. Rozważaną alternatywą dla OpenGL było inne API grupy Khronos: Vulkan. Vulkan został użyty m.in.\ w 3DGS.cpp, oferując jawną kontrolę nad pamięcią oraz dostęp do operacji na poziomie podgrup. To API jednak wymaga bardzo złożonej i rozbudowanej warstwy inicjalizacyjnej, co pochłonęłoby o wiele więcej czasu niż implementacja w OpenGL. OpenGL zatem został wybrany, gdyż jest wspierany przez sterowniki wszystkich głównych producentów GPU, jest wieloplatformowy (aczkolwiek w obecnej wersji API OpenGL jest uznawany za przestarzały na platformach Apple) i udostępnia wszystkie mechanizmy potrzebne do metody 3DGS.
\section{GLFW 3.4}
\label{glfw}
GLFW to lekka, otwartoźródłowa biblioteka napisana w C, odpowiadająca za okno aplikacji, kontekst w OpenGL czy obsługę urządzeń wejściowych~\cite{glfw34docs}. Jest wykorzystywana do ukrycia różnic między systemami operacyjnymi i ich systemami okienkowymi (Win32, X11, Wayland, Cocoa). Biblioteka zawiera również precyzyjny zegar systemowy użyty do pomiaru czasu klatki. Kolejną zaletą jest integracja z biblioteką Dear ImGui, dostarczona w postaci gotowego backendu. Bez niej zdarzenia trzeba byłoby przekazywać do interfejsu użytkownika ręcznie. Wersja 3.4 wprowadziła stabilne wsparcie dla protokołu Wayland, co ma znaczenie przy uruchamianiu aplikacji na współczesnych dystrybucjach Linuksa. Alternatywą była biblioteka SDL, która obejmuje szerszy zakres funkcji, obsługę dźwięku i kontrolerów. Ponieważ aplikacja nie wymaga warstwy audio, przewaga ta nie ma znaczenia, a mniejszy rozmiar przemawia na korzyść GLFW.
\section{GLAD 2}
\label{glad}
OpenGL definiuje jedynie interfejs funkcji, natomiast ich implementacja znajduje się w sterowniku karty graficznej i nie jest dostępna w postaci statycznie linkowanej biblioteki. Adresy funkcji muszą zostać pobrane za pośrednictwem mechanizmów systemowych, co wymuszałoby na programiście ręczne pozyskanie kilkuset wskaźników. GLAD tworzy loader na podstawie przekazanej wersji API oraz profilu~\cite{glad2repo}. Loader to kod źródłowy deklarujący i inicjalizujący wyłącznie te funkcje, które są potrzebne. W projekcie wykorzystano GLAD w wersji 2 generujący loader dla konfiguracji 4.3 bezpośrednio na etapie konfiguracji. Dzięki temu wygenerowany kod nie jest przechowywany w repozytorium, lecz powstaje automatycznie podczas budowania. Profil Core zapewnia korzyść w postaci wczesnego wykrywania błędów. Próba użycia funkcji usuniętej z profilu Core (co znaczy, że funkcja pochodziła z przestarzałego potoku stałofunkcyjnego) zakończy się błędem kompilacji. Alternatywą mogłaby być biblioteka GLEW, jednak ładuje ona pełen zestaw rozszerzeń niezależnie od potrzeb aplikacji. GLAD 2 ma również opcję generowania powtarzalnego wyniku, co jest ważne przy weryfikacji, że kolejne kompilacje projektu dają te same artefakty.
\section{GLM (OpenGL Mathematics)}
\label{glm}
GLM to biblioteka nagłówkowa napisana w C++, której typy mają wiernie odwzorowywać typy występujące w języku GLSL~\cite{glmrepo}. Biblioteka ta nie jest ograniczona tylko do funkcji znajdujących się w GLSL, ale jest rozszerzona o mechanizmy związane z przekształceniami macierzy, generowaniem liczb losowych i szumem. GLM przechowuje macierze w postaci kolumnowej, zgodnie z konwencją przyjętą w OpenGL, dzięki czemu można je przesyłać bez transpozycji. W projekcie biblioteka ta została użyta do wyznaczenia macierzy projekcji perspektywicznej, widoku, przekształceń kamery oraz do normalizacji wektorów. Biblioteka, będąc biblioteką nagłówkową, nie wymaga osobnego kompilowania ani linkowania, co przyspiesza integrację. Funkcje szablonowe są w pełni rozwijane w miejscu wywołania i podlegają optymalizacji kompilatora.
\section{Dear ImGui 1.91.8}
\label{imgui}
Dear ImGui to popularna biblioteka używana do tworzenia interfejsu użytkownika. Działa w tzw. immediate mode: kontrolki nie są trwałymi obiektami o własnym stanie, lecz są deklarowane od nowa w każdej klatce jako wywołania funkcji zwracające informację o interakcji użytkownika~\cite{imgui1918repo}. Dzięki takiemu podejściu znika konieczność synchronizowania stanu interfejsu ze stanem aplikacji. Jest to przydatne w aplikacji, w której parametry renderowania zmieniają się nieustannie. Biblioteka renderuje warstwę graficzną przy użyciu tego samego kontekstu OpenGL co scena, przez co nie wprowadza dodatkowych zależności systemowych i nie wpływa na wieloplatformowość aplikacji. Integrację zapewniają backendy GLFW i OpenGL 3, które przejmują obsługę zdarzeń wejściowych i wystawianie list rysowania. W tym projekcie ImGui odpowiada za panel sterowania parametrami renderowania, prezentację statystyk i przełączanie trybów renderowania. Interfejs jest kluczową przewagą nad narzędziem 3DGS.cpp, które udostępnia niewielką liczbę opcji konfiguracyjnych. Biblioteka jest mało wymagająca obliczeniowo, a jej narzut pamięciowy nie zależy od liczby splatów. Alternatywą dla Dear ImGui jest Qt, które oferuje bogatszy zestaw kontrolek, kosztem większej zależności zewnętrznej, bardziej złożonego procesu budowania i konieczności utrzymywania osobnej reprezentacji stanu interfejsu.
\section{ImPlot 1.0}
\label{implot}
ImPlot to biblioteka do wizualizacji wykresów stworzona na bazie Dear ImGui, dziedzicząca immediate mode oraz sposób integracji z aplikacją~\cite{implotrepo}. W projekcie odpowiada za warstwę analityczną. Za jej pomocą prezentowane są rozkłady nieprzezroczystości oraz skali Gaussianów, które są obecnie widoczne, a także histogram rozmiaru rzutowanych elips w przestrzeni ekranu. Dzięki tym informacjom można przeprowadzić analizę sceny -- na przykład rozpoznać nadmiar splatów o niskiej nieprzezroczystości, które spowalniają rendering. Wykresy reagują na ruch kamery i mogą służyć do analizy poszczególnych fragmentów sceny (co jest dodatkowo możliwe poprzez zapisanie ustawień kamery w pliku .json). ImPlot został jeszcze wykorzystany do zebrania pomiarów wydajnościowych przy różnych rodzajach sortowania splatów po głębokości. Rozwiązaniem alternatywnym byłby eksport danych do pliku i analiza ich w zewnętrznym narzędziu, co uniemożliwiałoby obserwację rozkładów w czasie rzeczywistym.
\section{Happly}
\label{happly}
W programie splat jest reprezentowany przez strukturę w kodzie. Sceny 3DGS są zapisywane w formacie .ply opracowanym na Uniwersytecie Stanforda. Format jest zbudowany z nagłówka deklarującego tekstowo elementy, ich właściwości i typy, jak i z danych właściwych, zapisanych w postaci tekstowej albo binarnej. W przypadku 3DGS na element vertex składa się kilkadziesiąt właściwości, takich jak pozycja, skala, kwaternion rotacji, nieprzezroczystość oraz współczynniki harmonik sferycznych. Ze względu na dużą liczbę rekordów często te pliki mają rozmiar setek megabajtów. Happly to nagłówkowa biblioteka udostępniająca ten format w postaci prostego interfejsu opartego na pobieraniu nazwanych parametrów wskazanego elementu~\cite{happlyrepo}. W projekcie została opakowana funkcją pomocniczą sprawdzającą, czy dana własność występuje. Dzięki temu można obsłużyć pliki pochodzące z różnych implementacji treningu lub różniące się liczbą zapisanych harmonik sferycznych. Cechą, która na to pozwala, jest samoopisywalność formatu PLY: z nagłówka można odczytać, jakie właściwości są w pliku. Dane odczytane z pliku zostają przetworzone przed zapisaniem ich do struktury splatu (np.\ funkcja sigmoidalna na nieprzezroczystości, funkcja wykładnicza na skali, normalizacja kwaternionu). Ta biblioteka została wybrana ze względu na łatwą integrację i brak zależności zewnętrznych. Alternatywą była biblioteka tinyply, która oferowała zbliżoną funkcjonalność, ale o bardziej rozbudowanym interfejsie, opartym na jawnym zarządzaniu buforami.
\section{Biblioteki pomocnicze}
\label{help_lib}
\subsection{stb}
\label{stb}
stb to zbiór jednoplikowych bibliotek nagłówkowych upowszechniony przez Seana Barretta. W projekcie wykorzystano dwa moduły: stb\_image\_write do zapisu zrzutu ekranu do pliku PNG oraz stb\_image do wczytywania obrazów (wykorzystywane przy prototypowaniu)~\cite{stbrepo}. Zaletą bibliotek stb jest brak jakichkolwiek zależności i licencja nienakładająca ograniczeń, a koszt integracji sprowadza się do dołączenia do projektu pojedynczego pliku nagłówkowego.
\subsection{nlohmann/json 3.11.3}
\label{json}
Biblioteka ta pozwala na obsługę formatu JSON w postaci interfejsu zbliżonego do kontenerów biblioteki standardowej~\cite{nlohmannJsonDocs}. W projekcie służy do zapisu i odczytu ustawień kamery oraz ścieżek przejazdów w testach wydajnościowych. Dzięki formatowi tekstowemu pliki konfiguracyjne pozostają czytelne i edytowalne poza aplikacją, co ułatwia ich przygotowanie do powtarzalnych scenariuszy pomiarowych. Zapis ustawień umożliwia porównanie wydajnościowe pomiędzy kolejnymi iteracjami implementacji.
\subsection{nativefiledialog-extended 1.2.1}
\label{filedialog}
Biblioteka ta udostępnia interfejs do otwierania systemowych okien wyboru pliku, wywołując pod spodem natywne mechanizmy danej platformy~\cite{nativefiledialogRepo}. Dzięki temu użytkownik korzysta z okna dialogowego występującego w jego środowisku, zamiast z implementacji przygotowanej wewnątrz aplikacji.

\section{CMake}
\label{cmake}
CMake to wieloplatformowe narzędzie służące do zarządzania procesem budowania aplikacji, które na podstawie opisu projektu generuje pliki dla natywnego narzędzia budującego danej platformy. Rozwiązanie to jest standardem w projektach C++ dzięki przenośności, ponieważ pozwala utrzymać jeden opis projektu dla wielu platform jednocześnie. Wymagana wersja CMake to 3.21, a dzięki FetchContent wszystkie zależności zewnętrzne są pobierane automatycznie poprzez klonowanie wskazanych repozytoriów~\cite{cmake320fetchcontent}. Dla każdej zależności została zadeklarowana konkretna wersja, co sprawia, że proces budowania jest powtarzalny i odporny na zmiany wprowadzone w repozytoriach źródłowych. Dzięki takiemu podejściu nie ma potrzeby przechowywania kodu bibliotek ani ręcznego instalowania wymaganych zależności. CMake pozwala również na kopiowanie po zakończonej kompilacji plików takich jak shadery i zasoby do katalogu wynikowego, dzięki czemu nie trzeba dodatkowo przygotowywać środowiska.